% ========================================================
% BÀI TẬP TOÁN 5 - KẾT NỐI TRI THỨC VỚI CUỘC SỐNG
% Bài 64: Biểu đồ hình quạt tròn
% Định dạng: BTTeX Part III - Trắc nghiệm điền số ngắn
% ========================================================

% Câu 1
\begin{ex}
  Một biểu đồ hình quạt tròn biểu diễn thể loại phim yêu thích của 100 học sinh: Hoạt hình $50\%$, Ca nhạc $25\%$, Thể thao $15\%$, Phim tài liệu $10\%$. Có bao nhiêu học sinh thích phim Hoạt hình?
  \shortans{50}
  \loigiai{
    $100 \times 50\% = 50$ học sinh.
  }
\end{ex}

% Câu 2
\begin{ex}
  Ở biểu đồ trên, có bao nhiêu học sinh thích phim Thể thao?
  \shortans{15}
  \loigiai{
    $100 \times 15\% = 15$ học sinh.
  }
\end{ex}

% Câu 3
\begin{ex}
  Ở biểu đồ trên, số học sinh thích phim Ca nhạc nhiều hơn phim Tài liệu là bao nhiêu bạn?
  \shortans{15}
  \loigiai{
    $25 - 10 = 15$ bạn.
  }
\end{ex}

% Câu 4
\begin{ex}
  Biểu đồ hình quạt tròn thể hiện diện tích trồng cây: Cây ăn quả chiếm $40\%$, Cây hoa chiếm $35\%$, còn lại là cây bóng mát. Tỉ số phần trăm cây bóng mát là bao nhiêu phần trăm?
  \shortans{25}
  \loigiai{
    $100\% - 40\% - 35\% = 25\%$.
  }
\end{ex}

% Câu 5
\begin{ex}
  Ở khu vườn trên rộng $2\,000\text{ m}^2$. Diện tích trồng cây ăn quả là bao nhiêu mét vuông?
  \shortans{800}
  \loigiai{
    $2\,000 \times 40\% = 800\text{ m}^2$.
  }
\end{ex}

% Câu 6
\begin{ex}
  Ở khu vườn trên, diện tích trồng hoa là bao nhiêu mét vuông?
  \shortans{700}
  \loigiai{
    $2\,000 \times 35\% = 700\text{ m}^2$.
  }
\end{ex}

% Câu 7
\begin{ex}
  Biểu đồ quạt tròn biểu diễn kết quả học tập của 200 học sinh: Loại Giỏi $30\%$, Khá $50\%$, Đạt $20\%$. Có bao nhiêu học sinh xếp loại Giỏi?
  \shortans{60}
  \loigiai{
    $200 \times 30\% = 60$ học sinh.
  }
\end{ex}

% Câu 8
\begin{ex}
  Ở trường trên, có bao nhiêu học sinh xếp loại Khá?
  \shortans{100}
  \loigiai{
    $200 \times 50\% = 100$ học sinh.
  }
\end{ex}

% Câu 9
\begin{ex}
  Biểu đồ quạt tròn biểu diễn cơ cấu các loại sách trong thư viện có 1 000 cuốn sách: Sách giáo khoa chiếm $60\%$, Truyện tranh chiếm $25\%$, Sách tham khảo chiếm $15\%$. Số truyện tranh là bao nhiêu quyển?
  \shortans{250}
  \loigiai{
    $1\,000 \times 25\% = 250$ quyển.
  }
\end{ex}

% Câu 10
\begin{ex}
  Ở thư viện trên, số sách giáo khoa là bao nhiêu quyển?
  \shortans{600}
  \loigiai{
    $1\,000 \times 60\% = 600$ quyển.
  }
\end{ex}

% Câu 11
\begin{ex}
  Một hình tròn biểu thị toàn bộ dữ liệu ứng với bao nhiêu phần trăm?
  \shortans{100}
  \loigiai{
    Toàn bộ hình tròn biểu thị 100%.
  }
\end{ex}

% Câu 12
\begin{ex}
  Biểu đồ quạt tròn chia làm 4 phần bằng nhau. Mỗi phần biểu diễn bao nhiêu phần trăm?
  \shortans{25}
  \loigiai{
    $100\% : 4 = 25\%$.
  }
\end{ex}

% Câu 13
\begin{ex}
  Biểu đồ quạt tròn có một nửa hình tròn biểu thị xe đạp. Tỉ số phần trăm của xe đạp là bao nhiêu phần trăm?
  \shortans{50}
  \loigiai{
    Một nửa là $50\%$.
  }
\end{ex}

% Câu 14
\begin{ex}
  Một đàn gia cầm có 500 con. Biểu đồ quạt cho biết: Gà chiếm $50\%$, Vịt chiếm $30\%$, Ngan chiếm $20\%$. Số con ngan là bao nhiêu con?
  \shortans{100}
  \loigiai{
    $500 \times 20\% = 100$ con.
  }
\end{ex}

% Câu 15
\begin{ex}
  Ở đàn gia cầm trên, số con vịt là bao nhiêu con?
  \shortans{150}
  \loigiai{
    $500 \times 30\% = 150$ con.
  }
\end{ex}

% Câu 16
\begin{ex}
  Ở đàn gia cầm trên, số con gà nhiều hơn số con vịt bao nhiêu con?
  \shortans{100}
  \loigiai{
    Gà: $500 \times 50\% = 250$. Hơn: $250 - 150 = 100$ con.
  }
\end{ex}

% Câu 17
\begin{ex}
  Khảo sát phương tiện đến trường của 120 học sinh: Đi bộ $20\%$, Xe đạp $50\%$, Bố mẹ chở $30\%$. Số bạn tự đi xe đạp là bao nhiêu bạn?
  \shortans{60}
  \loigiai{
    $120 \times 50\% = 60$ bạn.
  }
\end{ex}

% Câu 18
\begin{ex}
  Ở khảo sát trên, số bạn đi bộ đến trường là bao nhiêu bạn?
  \shortans{24}
  \loigiai{
    $120 \times 20\% = 24$ bạn.
  }
\end{ex}

% Câu 19
\begin{ex}
  Biểu đồ quạt phân loại rác thải của một xóm: Rác hữu cơ $60\%$, Rác tái chế $30\%$, Rác nguy hại $10\%$. Tổng lượng rác là $200\text{ kg}$. Lượng rác tái chế là bao nhiêu ki-lô-gam?
  \shortans{60}
  \loigiai{
    $200 \times 30\% = 60\text{ kg}$.
  }
\end{ex}

% Câu 20
\begin{ex}
  Ở xóm trên, lượng rác hữu cơ là bao nhiêu ki-lô-gam?
  \shortans{120}
  \loigiai{
    $200 \times 60\% = 120\text{ kg}$.
  }
\end{ex}

% Câu 21
\begin{ex}
  Một tháng có chi tiêu gia đình: Ăn uống $45\%$, Tiết kiệm $15\%$, Chi tiêu khác $40\%$. Nếu thu nhập là $20$ triệu đồng thì số tiền tiết kiệm được là bao nhiêu triệu đồng?
  \shortans{3}
  \loigiai{
    $20 \times 15\% = 3$ triệu đồng.
  }
\end{ex}

% Câu 22
\begin{ex}
  Ở bài trên, số tiền dành cho ăn uống là bao nhiêu triệu đồng?
  \shortans{9}
  \loigiai{
    $20 \times 45\% = 9$ triệu đồng.
  }
\end{ex}

% Câu 23
\begin{ex}
  Một trang trại có 800 cây ăn quả gồm: Xoài $35\%$, Nhãn $40\%$, Vải $25\%$. Số cây nhãn là bao nhiêu cây?
  \shortans{320}
  \loigiai{
    $800 \times 40\% = 320$ cây.
  }
\end{ex}

% Câu 24
\begin{ex}
  Ở trang trại trên, số cây vải là bao nhiêu cây?
  \shortans{200}
  \loigiai{
    $800 \times 25\% = 200$ cây.
  }
\end{ex}

% Câu 25
\begin{ex}
  Một hình quạt có số đo góc ở tâm là $90^\circ$. Hình quạt này chiếm bao nhiêu phần trăm của hình tròn?
  \shortans{25}
  \loigiai{
    $\frac{90}{360} = \frac{1}{4} = 25\%$.
  }
\end{ex}
